\begin{abstract}
La radiometría es la rama de la óptica que estudia la propagación y distribución energética de la luz en función de magnitudes como flujo radiante, irradiancia, intensidad radiante y radiancia. En este experimento, se estudió la relación entre la irradiancia y la distancia a la fuente luminosa para dos tipos de emisores: un LED y un láser. Se realizaron mediciones variando la distancia entre el sensor de irradiancia y la fuente en incrementos de ($5 \pm 2$)mm, ($10 \pm 2$)mm y ($50 \pm 2$)mm. Se compararon los resultados con dos modelos teóricos: la ley del inverso del cuadrado, que predice que la irradiancia de una fuente perfectamente puntual disminuye con \( r^{-2} \), y el modelo de haz colimado, en el que la irradiancia es constante. Los resultados experimentales mostraron que el LED sigue la tendencia de la ley del inverso del cuadrado con un exponente de decaimiento \( m = -1.94 \pm 0.00647 \), mientras que el láser presenta un exponente significativamente menor, \( m = -0.0182 \pm 0.00167 \), indicando una propagación casi colimada (i.e., de divergencia tendiendo a cero para distancias de hasta 400mm acorde a este estudio). Estos resultados validan la dependencia de la irradiancia con la naturaleza de la fuente y confirman la aplicabilidad de los modelos radiométricos en la caracterización de fuentes de luz.
\end{abstract}